\section{Related Work}
\label{sec:related}

% R1: memory systems (A/B: clearly better than LongMemEval's memory-methods paragraph; positioning fixes applied)
\paragraph{Memory systems for long conversations.}
Memory systems for chat assistants differ in when they organize what they store. MemGPT treats the context window as
virtual memory and lets the agent move and edit content between it and external storage~\citep{packer2023memgpt}.
Most later systems instead store an LLM-processed version of the history and reorganize it at write time, as new turns
or sessions arrive. Mem0 extracts facts and decides for each new fact whether to add it, update an existing one, or
delete one~\citep{chhikara2025mem0}. A-Mem links each new note to related notes and updates the tags and context of
existing ones~\citep{xu2025amem}. MemoryOS moves content between short-, mid-, and long-term
stores~\citep{kang2025memoryos}. Zep builds a temporal knowledge graph whose edges are invalidated rather than deleted
when facts change~\citep{rasmussen2025zep}. Hindsight organizes memory into four networks of world facts, agent
experiences, synthesized entity summaries, and evolving beliefs, and updates them as conversations
arrive~\citep{latimer2025hindsight}. Other systems defer or limit this work: LightMem consolidates long-term memory
offline~\citep{fang2025lightmem}, LycheeMemory encodes each finalized segment of dialogue into context-independent
typed records rather than merging it into existing memory~\citep{li2026lychee}, and LeanMem updates only the event
memories that evolve while keeping source-grounded records unchanged~\citep{liao2026leanmem}. Even these decide what to
merge before any question is asked. A related line handles conflicting memories when they are written: MOSAIC checks
each new fact against its graph neighbors and updates or deletes contradicted entries~\citep{zhao2026mosaic}, TOKI
types contradiction resolution as bitemporal operators~\citep{wang2026toki}, and TEPA revokes stale entries while
keeping them for audit~\citep{zhou2026tepa}. \method{} keeps one write-time step, segment-level extraction of atomic records as
in LycheeMemory, and defers everything that relates records to one another (merging, resolving conflicts, and
summarizing) to read time, where it is conditioned on the question. The price is extra computation for every
question, which we measure in \cref{sec:cost}.

% R2: reading strategies and query-focused compression
\paragraph{Organizing retrieved context for the reader.}
A second line of work organizes text for the reader, either when the index is built or when a query arrives.
RAPTOR recursively clusters and summarizes chunks into a tree that is queried at several levels of
abstraction~\citep{sarthi2024raptor}, and GraphRAG summarizes communities of an entity graph and answers global
questions by map-reduce over these summaries~\citep{edge2024graphrag}; both organize at index time. At query time,
Chain-of-Note trains a model to write a reading note on each retrieved document before answering~\citep{yu2023chainofnote},
RECOMP trains compressors that shorten retrieved documents into extractive or abstractive summaries~\citep{xu2024recomp},
ReadAgent prompts a model to compress pages into gists and look pages up again when needed~\citep{lee2024readagent}, and
Astute RAG consolidates retrieved and internal knowledge and resolves conflicts between them~\citep{wang2024astute}.
Self-RAG~\citep{asai2024selfrag} and Adaptive-RAG~\citep{jeong2024adaptive} decide when and how much to retrieve, which
is complementary to organizing what has been retrieved. Closest to us, LongMemEval finds that Chain-of-Note-style reading
and a structured presentation of retrieved sessions improve accuracy on long chat histories~\citep{wu2025longmemeval}.
Two recent memory systems also move work to query time. Pull keeps raw turns behind a deterministic index and lets the
reader expand only the turns it needs~\citep{chen2026pull}; it does not rewrite what it retrieves. DeferMem is the
closest to our method: it retrieves a broad, noisy candidate set from raw history at query time and distills it into
query-specific evidence with a model trained by reinforcement learning~\citep{yin2026defermem}, evaluated on LoCoMo
and LongMemEval. \method{} differs in three ways: it needs no training, it compiles from extracted records rather
than raw turns and passes the records to the reader along with the notes, and its notes have a fixed structure for
memory-specific failures (updates, conflicts, gaps in coverage, standing instructions). We also evaluate at the 1M
tier of BEAM, where DeferMem has not been tested.
We build on these observations rather than claim reader-side notes as new. Our compilation step is one prompted call per
question, and its format is specific to conversational memory: it orders evidence by time and marks facts that were
later updated or contradicted, parts of the question the memory does not cover, and the user's standing instructions.
Zep tracks the validity of facts at write time; we do so at read time over storage that never changes. Our main
contribution, however, is not the format but a controlled comparison of organizing at write time and at read time,
with retrieval, prompt, and judge held fixed, across readers that differ in how well they use long contexts.

% R3: how readers use long contexts (A/B: clearly better than Lost-in-the-Middle 6.2; edits applied)
\paragraph{How readers use long contexts.}
The \emph{reader}, the model that answers from what the memory retrieves, uses information in the middle of a long
input less reliably than information at its ends~\citep{liu2024lost}, and its effective context is often much shorter
than the window it accepts~\citep{hsieh2024ruler}. These findings suggest that giving the reader more memory context
need not raise accuracy. We observe this in memory systems as well, and we find that it varies sharply across
readers: the long-context gap that our compilation step (\cref{sec:method}) closes for one reader is already small
for another (\cref{sec:reader}).

% R4: evaluating memory systems
\paragraph{Evaluating long-term memory.}
LoCoMo~\citep{maharana2024locomo}, LongMemEval~\citep{wu2025longmemeval}, and BEAM~\citep{tavakoli2025beam} test
memory over conversations ranging from about 9K tokens (LoCoMo) to 10M tokens (BEAM). LongMemEval and BEAM grade
free-form answers with an LLM judge, whose verdicts depend on the judge model and its prompt~\citep{zheng2023judging}.
Numbers reported in different papers therefore differ in reader, answer prompt, context budget, and judge, and are
not directly comparable. RENDER fixes the conversation and varies only the form in which evidence reaches the reader,
and finds spreads of 25--49 points across forms on LongMemEval~\citep{si2026render}; \citet{jiang2026presentation}
find that most of the gain of a fine-grained revision ledger disappears once its presentation is matched by a flat
baseline; and MemConflict separates retrieval from use of retrieved evidence under conflicting
memories~\citep{tao2026memconflict}. Our audit makes the same kind of argument for budget and answer prompt on
BEAM, and adds a split of each system's lost points into storage, retrieval, and reading. We do not propose a new benchmark. Within each comparison we evaluate every system with the
same reader, prompt, and judge, run each condition three times, and report intervals clustered by conversation, so
that for a given reader the differences between rows reflect the memory rather than the evaluation setup; we then
repeat the main comparison with three readers (\cref{sec:setting}).
